% !PRINTED 772
\weberpage{772}

%<<CONT
integer $x$ passes into an integer which, multiplied by $p$, cannot be $= 1$.

With this the irreducibility of $f_n(x)$ is proved in general.

\begin{center}
\rule{4cm}{0.4pt}
\end{center}

Now that the irreducibility of the cyclotomic equation
$f_n(x) = 0$ has been established in general, it follows, as in \S.\,134 of the
first volume, that it also remains irreducible when an arbitrary root of unity is adjoined whose degree is relatively
prime to $n$.

\gesperrt{Kronecker} takes, in a certain sense, the opposite path\footnote{Mém. sur les facteurs irréd. de l'expression $(x^n-1)$. \gesperrt{Liouville}'s
Journal, Bd. 19 (1854). A simple proof of the irreducibility of the general cyclotomic equation, resting on the theory of the higher congruences,
has been given by \gesperrt{Dedekind} [\gesperrt{Crelle}'s Journal für Mathematik, Bd. 58 (1859)].}.

He first proves, under the assumption that $n$ is a
power of a prime $p$, the irreducibility of $f_n(x)$ after
the adjunction of a root of unity whose degree is not divisible by $p$,
and from this he derives the general proof of irreducibility, as we shall now show briefly.

\begin{center}
\rule{4cm}{0.4pt}
\end{center}

Let now
\[
n = p^k
\]
be a power of a prime $p$, and let
\weq{8}{f_n(x) = x^{p^{k-1}(p-1)} + x^{p^{k-1}(p-2)} + \cdots + x^{p^{k-1}} + 1}
be the function whose roots are the primitive $n$th roots of unity.
Let further $\alpha$ be any root of unity whose degree $m$
is not divisible by $p$, and let $\nu$ be the degree of the field $R(\alpha)$.
\gesperrt{It is to be proved that $f_n(x)$ is irreducible in the field $R(\alpha)$.}

Every number of the field $R(\alpha)$ can be represented in one and only
one way in the form:
\weq{9}{\varphi(\alpha) = m_0 + m_1 \alpha + \cdots + m_{\nu-1} \alpha^{\nu-1},}
where $m_0, m_1, \dots, m_{\nu-1}$ are rational numerical coefficients.

If $\chi(x) = 0$ is the rational equation of lowest degree
having $\alpha$ as its root, and
\weq{10}{\chi(x) = x^\nu + c_1 x^{\nu-1} + \cdots + c_{\nu-1} x + c_\nu,}
%CONT>>
